\subsection{Dataset and Evaluation Protocol}
\label{sec:dataset_and_metrics}

\paragraph{Dataset and Partitioning.}
We evaluated the detector on the KITTI 3D Object Detection dataset~\cite{geiger2012kitti}. The dataset contains $7{,}481$ calibrated LiDAR scans with synchronized rectified camera images and 3D bounding box annotations from road scenes. We used the 50/50 split introduced by Chen et al.~\cite{chen20153d}, with $3{,}712$ training frames and $3{,}769$ validation frames. Ground-truth sampling databases were generated exclusively from the $3{,}712$ training frames to prevent validation data leakage.

\paragraph{Spatial Region of Interest (ROI) and Target Parameterization.}
Following Section~\ref{sec:problem_formulation}, we aligned the LiDAR coordinate frame with the vehicle body (forward $X$, leftward $Y$, upward $Z$). We restricted perception to the Cartesian Region of Interest:
\begin{equation}
    \begin{aligned}
        x &\in [0.0, 70.4]\,\text{m}, \\
        y &\in [-40.0, 40.0]\,\text{m}, \\
        z &\in [-2.5, 1.0]\,\text{m}.
    \end{aligned}
\end{equation}
We applied an open-boundary inward margin of $\epsilon = 0.001\,\text{m}$ along all three axes. A horizontal grid resolution of $\Delta_x = \Delta_y = 0.1\,\text{m}$ produced an input tensor of $H_{\mathrm{in}} \times W_{\mathrm{in}} = 800 \times 704$ cells. An output stride of $s_{\mathrm{out}} = 4$ gave feature maps of size $H_{\mathrm{out}} \times W_{\mathrm{out}} = 200 \times 176$. We evaluated Car, Pedestrian, and Cyclist detections.

\paragraph{Evaluation Metrics.}
We evaluated oriented 2D Bird's-Eye View (BEV) footprints on the KITTI validation set using 40-point interpolated Average Precision (AP $R_{40}$). This metric averages interpolated precision at 40 recall positions $\mathcal{R}_{40} = \{1/40, 2/40, \dots, 1.0\}$:
\begin{equation}
    \mathrm{AP}|_{R_{40}} = \frac{1}{40} \sum_{r \in \mathcal{R}_{40}} \max_{r' \ge r} p(r').
\end{equation}
We used nominal oriented BEV Intersection-over-Union (IoU) thresholds of $\tau_{\mathrm{IoU}} = 0.70$ for Car and $\tau_{\mathrm{IoU}} = 0.50$ for Pedestrian and Cyclist. We grouped detections into three difficulty regimes:
\begin{itemize}
    \item \textbf{Easy}: 2D bounding box height $\ge 40\,\text{px}$, maximum occlusion level = 0 (fully visible), maximum truncation $\le 15\%$.
    \item \textbf{Moderate}: 2D bounding box height $\ge 25\,\text{px}$, maximum occlusion level = 1 (partially occluded), maximum truncation $\le 30\%$.
    \item \textbf{Hard}: 2D bounding box height $\ge 25\,\text{px}$, maximum occlusion level = 2 (difficult to see), maximum truncation $\le 50\%$.
\end{itemize}
Following autonomous driving literature, \textbf{Moderate macro mAP} ($\mathrm{mAP}_{\mathrm{M}}$) was adopted as the primary accuracy summary metric:
\begin{equation}
    \mathrm{mAP}_{\mathrm{M}} = \frac{1}{3} \sum_{c \in \{\mathrm{Car}, \mathrm{Ped}, \mathrm{Cyc}\}} \mathrm{AP}_{c, \mathrm{M}}.
\end{equation}
We also reported the unweighted mean across all nine class-difficulty combinations:
\begin{equation}
    \mathrm{mAP}_{9} = \frac{1}{9} \sum_{c \in \mathcal{C}} \sum_{d \in \{\mathrm{E}, \mathrm{M}, \mathrm{H}\}} \mathrm{AP}_{c, d}.
\end{equation}
E, M, and H denote Easy, Moderate, and Hard. $\mathrm{mAP}_{9}$ averages the nine class--difficulty AP values; it does not denote AP at an IoU threshold of 0.9.

\paragraph{Unified Benchmark Protocol and Evaluation Environment.}
To ensure strict methodological fairness and eliminate discrepancies stemming from disparate evaluation implementations or hardware environments, all evaluated detectors in our benchmark comparison are evaluated under a unified, loader-aligned evaluation protocol. The evaluation code strictly adheres to the official KITTI devkit evaluation protocol (\texttt{kitti\_eval} protocol standard in OpenPCDet and mmdetection3d), computing the official 40-point interpolated Average Precision ($R_{40}$) metric across all nine class--difficulty evaluation cells on identical validation scenes ($3{,}769$ frames) using standard IoU thresholds ($0.70$ for Car, $0.50$ for Pedestrian and Cyclist). All latency and throughput measurements are benchmarked on the same dedicated NVIDIA Tesla T4 GPU platform (16\,GB GDDR6 VRAM, CUDA 13.0, PyTorch 2.11.0) under identical single-batch conditions ($\text{batch size} = 1$) and FP32 numerical precision.
\subsection{Implementation Details}
\label{sec:implementation_details}

\begin{table}[t]
\centering\small
\caption{Experimental settings, architecture hyperparameters, and training configuration for MobilePIXORNeXt. Training uses a deterministic single-seed setup ($\text{seed} = 42$) with PyTorch 2.11.0 and CUDA 13.0. Timing uses a dedicated NVIDIA Tesla T4 GPU (16\,GB GDDR6 VRAM), batch size 1, and FP32 precision.}
\label{tab:setup}
\begin{tabular}{@{}>{\raggedright\arraybackslash}p{.38\columnwidth}>{\raggedright\arraybackslash}p{.56\columnwidth}@{}}
\toprule
\textbf{Configuration Element} & \textbf{Specification / Value} \\
\midrule
LiDAR Input Representation & Strata-BEV14 (14 channels; NumPy/Numba CPU) \\
Spatial ROI bounds ($X, Y, Z$) & $[0.00, 70.40]\,\text{m} \times [-40.00, 40.00]\,\text{m} \times [-2.50, 1.00]\,\text{m}$ \\
Grid resolution ($\Delta_x, \Delta_y, \Delta_z$) & $0.10\,\text{m} \times 0.10\,\text{m} \times 0.10\,\text{m}$ \\
Input / Output Grid Dimensions & $800 \times 704$ / $200 \times 176$ ($s_{\mathrm{out}} = 4$) \\
Detector Parameters & 753,569 deployable (753,581 trainable with 12 T-SBUW logits); 9.34\,MB FP32 \\
Backbone Stages ($C_2, C_3, C_4$) & Widths: (48, 96, 128); Depths: (3, 4, 2) \\
Inverted Bottleneck Expansion & $e = 2.50$; $7 \times 7$ Depthwise Convolution \\
Attention Mechanism & Parameter-free Block-SimAM (9 residual blocks) \\
Multiscale Neck & Scale-Gated FPN (SG-FPN, $F_{32}$ fusion width) \\
Prediction Towers & GroupedHeader: Car ($K_1=1$), Ped/Cyc ($K_2=2$) \\
Loss Formulation & OGA (MGIoU + $\pi$-Corner) + IQA + T-SBUW \\
Data Augmentation & Physics-Consistent GT (PC-GTA) + Global \\
Optimizer / Base Learning Rate & AdamW (fused) / $\eta = 7.00 \times 10^{-4}$ \\
Minimum LR / Training Schedule & $1.00 \times 10^{-6}$ / 100 epochs: 10 warmup + 90 cosine \\
Final Reported Checkpoint & Converged deployment checkpoint evaluated at epoch 100 \\
Weight Decay & $\lambda = 1.00 \times 10^{-3}$ (2D+ tensors only; biases/norm: 0.00) \\
Gradient Clipping Norm & Maximum $L_2$ norm = 10.00 \\
Batch Size (Train / Val) & 16 / 32 \\
Numerical Precision & Automatic Mixed Precision (AMP / FP32) \\
Hardware Environment & NVIDIA Tesla T4 GPU (16\,GB GDDR6, 98.44\,MiB peak) \\
\bottomrule
\end{tabular}
\end{table}

Table~\ref{tab:setup} summarizes the primary implementation settings and architectural hyperparameters across all pipeline components.

\paragraph{Network and Model Configuration.}
MobilePIXORNeXt processes 14-channel Strata-BEV14 tensors. We used a two-layer $3 \times 3$ convolutional stem ($32$ channels, stride 2) and three backbone stages with channel widths $(48, 96, 128)$ and block depths $(3, 4, 2)$. The stages used \texttt{MobilePIXORNeXt-S} blocks with $7 \times 7$ depthwise convolutions, an inverted expansion ratio of $e=2.5$, SiLU activations, and LayerScale initialized to $10^{-5}$. Parameter-free Block-SimAM attention was embedded in all nine residual blocks. The Scale-Gated Feature Pyramid Network (SG-FPN) fused stages C4, C3, and C2 into a $F_{32}$ (32-channel) representation at stride 4 ($200 \times 176$). Decoupled GroupedHeader towers processed Car ($K_1=1$) and Pedestrian/Cyclist ($K_2=2$) independently, with an initial $3 \times 3$ Conv-BN-SiLU transformation per group. The detector contained $753{,}569$ deployable parameters ($753{,}581$ trainable parameters during training including 12 transient T-SBUW multi-task logits), with an FP32 checkpoint size of $9.34\,\text{MB}$.

\paragraph{Training Schedule and Optimization.}
We optimized the model with AdamW using fused CUDA kernels. The base learning rate was $\eta_0 = 7 \times 10^{-4}$, with a minimum of $\eta_{\min} = 1.0 \times 10^{-6}$. Weight decay of $\lambda = 0.001$ applied only to multidimensional weight matrices ($\ge 2\mathrm{D}$ tensors). Biases, Batch Normalization affine parameters, and task-uncertainty logits were exempt ($\lambda = 0.0$). Gradients were clipped to a maximum $L_2$ norm of $10.0$. The standard training schedule covered 100 epochs, ending at epoch 100, comprising 10 linear warmup epochs and 90 cosine annealing epochs down to $\eta_{\min}$. To isolate and quantify each architectural component under strictly identical computational budgets, all exploratory ablation experiments (A1--A10) were trained under a standardized 50-epoch schedule (4 linear warmup epochs and 46 cosine annealing epochs), and evaluated at their 50-epoch checkpoints. Following this ablation study, the resulting optimal architecture (A10, MobilePIXORNeXt) and the baseline detector (A0, MobilePIXOR) were trained for a full 100-epoch schedule (10 linear warmup epochs and 90 cosine annealing epochs down to $\eta_{\min}$) to achieve complete asymptotic convergence for the definitive benchmark comparisons against external state-of-the-art detectors (matching the 100-epoch training budget of the multi-seed robustness evaluations under seeds 43 and 44). Section~\ref{sec:ablations} validates schedule convergence by comparing the 100-epoch deployment model against the 50-epoch ablation checkpoint. Temperature-Softmax Bounded Uncertainty Weighting (T-SBUW, Section~\ref{sec:losses_and_oga}) balanced the task losses. Section~\ref{sec:ablations} evaluates its empirical convergence behavior.

\paragraph{Data Augmentation Pipeline.}
Training used Physics-Consistent Ground-Truth Augmentation (\texttt{PC-GTA}; Section~\ref{sec:pc_gta}). The final reported run enabled local ground fitting, static obstacle clearance, range-adaptive point subsampling, and visibility protection. The static obstacle clearance threshold was set to $N_{\mathrm{thresh}} = 2$ points within the expanded footprint, and estimated ground elevation was bounded within $[-2.5, -0.8]\,\text{m}$. LOS rejection, shadow masking, and radiometric calibration were available as configurable options but disabled in the primary benchmark. Ground height was estimated within $R(d) = \operatorname{clip}(2.5 + 0.03\max(0, d - 15), 2.5, 4.0)\,\text{m}$. GT sampling was applied with probability $0.5$. Subsequent world transforms included flipping along the $X$-axis with probability $0.5$, rotation $\theta \sim \mathcal{U}(-45^\circ,45^\circ)$, and isotropic scaling $s \sim \mathcal{U}(0.95,1.05)$. World translation was applied with probability $0.5$, using independent Gaussian noise with standard deviations $(0.2,0.2,0.1)\,\text{m}$ along $(X,Y,Z)$.

\begin{table*}[t]
\centering\small
\caption{Quantitative comparison of Bird's-Eye View (BEV) object detection performance on the KITTI validation benchmark. Average Precision (AP, \%) is evaluated under the 40-point interpolation protocol ($R_{40}$) at Moderate difficulty ($\text{IoU} \ge 0.70$ for Car; $\text{IoU} \ge 0.50$ for Pedestrian and Cyclist). All experiments are conducted using the unified evaluation implementation and benchmarked on the same NVIDIA Tesla T4 GPU platform. The top partition benchmarks single-stage real-time and edge-oriented LiDAR detectors (spanning ultra-compact BEV representations to pillar and 3D sparse voxel networks); the bottom partition provides external reference with heavyweight two-stage 3D anchor-based detectors. For single-stage methods, the best results are highlighted in \textbf{bold}, and the second-best results are \underline{underlined}.}
\label{tab:comparison}
\setlength{\tabcolsep}{3pt}
\begin{tabular}{@{}lrrrrrr@{}}
\toprule
\textbf{Method} & \textbf{\shortstack{Params\\(M) $\downarrow$}} & \textbf{\shortstack{Car AP$_{\mathrm{M}}$\\(\%)}} & \textbf{\shortstack{Pedestrian\\AP$_{\mathrm{M}}$ (\%)}} & \textbf{\shortstack{Cyclist\\AP$_{\mathrm{M}}$ (\%)}} & \textbf{\shortstack{mAP$_{\mathrm{M}}$\\(\%)}} & \textbf{\shortstack{mAP$_{9}$\\(\%)}} \\
\midrule
PIXOR (Car only)~\cite{yang2018pixor} & 2.135 & 72.02 & -- & -- & -- & -- \\
MobilePIXOR~\cite{nguyen2024indoor} & \textbf{0.598} & 75.52 & 28.16 & 20.53 & 41.40 & 42.27 \\
PointPillars~\cite{lang2019pointpillars} & 4.835 & \underline{88.02} & 56.01 & 66.34 & 70.13 & 72.25 \\
SECOND~\cite{yan2018second} & 5.326 & \textbf{88.56} & \underline{56.56} & \textbf{71.16} & \textbf{72.09} & \textbf{73.79} \\
\textbf{MobilePIXORNeXt (Ours)} & \underline{0.754} & 86.90 & \textbf{60.78} & \underline{66.81} & \underline{71.50} & \underline{73.45} \\
\midrule
\multicolumn{7}{@{}l}{\textit{Heavyweight 3D Anchor-Based Reference Detectors}} \\
PointRCNN~\cite{shi2019pointrcnn} & 4.039 & 88.42 & 57.84 & 73.87 & 73.38 & 75.31 \\
Part-A$^2$-Anchor~\cite{shi2020part} & 63.815 & 90.05 & 64.20 & 74.60 & 76.28 & 78.04 \\
PV-RCNN~\cite{shi2020pvrcnn} & 13.116 & 90.29 & 58.48 & 74.54 & 74.44 & 76.49 \\
\bottomrule
\end{tabular}
\par\vspace{2pt}
{\footnotesize \textit{Note.} $\text{mAP}_{\mathrm{M}}$ represents the mean of the three Moderate class AP scores; $\text{mAP}_{9}$ averages all nine class--difficulty evaluation cells. Dashes (--) indicate unsupported classes for the single-class PIXOR detector. All comparative experiments are conducted under the unified evaluation file and executed on the same NVIDIA Tesla T4 GPU hardware platform with batch size 1 and FP32 precision to ensure a fair and controlled benchmark. Complete pipeline latencies and runtime component breakdowns are reported in Table~\ref{tab:efficiency}.}
\end{table*}


\paragraph{Batching, Precision, and Hardware Environment.}
Training used a physical batch size of 16 on a single GPU; validation used a batch size of 32. DataLoader processes used 6 persistent worker threads, pinned memory, and a prefetch factor of 2. Training was accelerated via Automatic Mixed Precision (AMP; utilizing Bfloat16 on Ampere/Ada cloud accelerators or FP16 on Turing). Loss evaluation, coordinate regression, and rotated NMS used full FP32 precision for numerical stability. Final MobilePIXORNeXt timing and accuracy evaluation used an NVIDIA Tesla T4 GPU (16\,GB GDDR6 VRAM, Turing architecture), an x86\_64 host CPU, Linux OS, and PyTorch 2.11.0 with CUDA 13.0. Peak allocated evaluation GPU memory was $98.44\,\text{MiB}$ ($103.22\,\text{MB}$ decimal). All evaluated detectors, including the MobilePIXOR baseline, are benchmarked on the same NVIDIA Tesla T4 GPU platform (detailed pipeline latency and throughput are analyzed in Section~\ref{sec:latency_analysis}). Primary comparative metrics and architectural ablations utilize deterministic evaluation ($\text{seed} = 42$) on the same validation split. In addition, multi-seed statistical robustness across distinct initializations ($\text{seeds} \in \{42, 43, 44\}$) is comprehensively analyzed in Table~\ref{tab:seed_robustness} (Section~\ref{sec:benchmark_comparison}).
